% Part of main.tex; compile main.tex.
\subsection{Discussion}
\label{sec:discussion}
In this section, we discuss when \texttt{sleep} is an acceptable stand-in for
computation, the practical implications of our results, and the threats to
their validity.

Our measurements show that \texttt{sleep} is a fair stand-in for computation
only under narrow conditions. The background engine must have a low CPU
demand, as POSIX AIO does. In addition, the computation must neither
communicate nor compete for the hardware thread of the engine, and the host
thread must actually be idle. Even under these conditions, the halo exchange
and the heat mini-app make the compute phase of POSIX AIO \SlowPosixHalo and
\SlowPosixHeat$\times$ longer. For engines that perform real work on the
hardware thread of the application, such as the HDF5 asynchronous VOL and
ADIOS2 in Guided mode, the application runs up to \ExcmaxpctHdf\% and
\ExcmaxpctAdios\% longer than \texttt{sleep} predicts. The same holds for
engines that need CPU time to progress, such as non-blocking MPI-IO, which
runs up to \ExcmaxpctMpiio\% longer. Moving the server of the DMS to a
reserved core is not sufficient either: PDC runs up to \ExcmaxpctPdc\%
longer, because its clients poll for progress on the cores of the
application. Furthermore, the advantage of
\texttt{sleep} persists across compute-to-I/O ratios and across output
sizes up to 480\,GiB per run.

Three practical observations follow from these results. First, the same
contention appears in different places for different engines. It appears in
the compute phase for POSIX, HDF5, and ADIOS2, and as exposed I/O for MPI-IO.
Reporting only the I/O time hides the former, and reporting only the compute
time hides the latter, while the end-to-end time captures both. Second,
process placement matters more for asynchronous I/O than the literature
suggests. Reserving hardware threads for the background engine recovers part
of the loss for HDF5 and turns non-blocking MPI-IO from a loss into a small
gain, but it costs compute resources. Third, GPU-resident computation is the
case in which \texttt{sleep} comes closest to reality. With pinned ranks, the
asynchronous run time with a GPU kernel matched the \texttt{sleep}
prediction within a few percent for POSIX, HDF5, and MPI-IO, while the same
ranks running DGEMM on the CPU ran up to \GpuexcpctHdfDgemm\% longer.
However, ADIOS2 still ran longer than predicted, and \texttt{sleep} still
inflated the speedups by slowing down the synchronous baseline.

\paragraph{Threats to validity.}
\label{sec:threats}
Our platform is a single node with local NVMe storage. It reproduces
contention for cores and memory bandwidth, but not contention for the
network links that MPI traffic and parallel file system traffic share. Our
halo surrogate approximates this contention only with intra-node
shared-memory transfers. In addition, we emulate a small client cache with a
cgroup memory limit rather than with the client of a real parallel file
system, and 4 ranks share one GPU in the GPU experiments, whereas production
GPU nodes typically have one GPU per rank.

Every configuration was run once, and the I/O time on this node varies
considerably between runs, with a median difference of \VarMedian\% and up
to \VarMax\% between two I/O-only runs (Section~\ref{sec:workload}). It is
important to note one systematic case. In the main experiment, the
synchronous run with \texttt{sleep} had the slowest synchronous I/O of all
surrogates for every backend: \SynciosleepPosix\,s per step for POSIX,
compared to \SynciootherminPosix to \SyncioothermaxPosix\,s for the other
surrogates, and \SynciosleepAdios\,s for ADIOS2, compared to
\SynciootherminAdios to \SyncioothermaxAdios\,s. Because this run was the
first one after the I/O-only baseline, we repeated the synchronous runs with
spin and \texttt{sleep} in the opposite order, with spin first. The result
did not change. The synchronous I/O after a \texttt{sleep} phase was again
slower in every backend, by \RecheckpctPosix\% for POSIX, \RecheckpctMpiio\%
for MPI-IO, \RecheckpctHdf\% for HDF5, and \RecheckpctAdios\% for ADIOS2.
Therefore, this is not an ordering effect. Instead, synchronous I/O appears
to run slower after the cores have been idle, possibly because of the power
states and clock frequency of idle cores, which we cannot inspect without
administrator access. This is a second way in which \texttt{sleep} distorts
an evaluation: besides hiding contention in the asynchronous run, it slows
down the synchronous baseline and thereby inflates the speedup. For this
reason, we
base our main comparison on asynchronous run times, which do not depend on
the synchronous runs (Section~\ref{sec:metrics}). Individual values of this
comparison are also uncertain, and effects of a few percent, such as
\ExcpctAdiosHeat\% for ADIOS2 with the heat mini-app, are within the
run-to-run variation. Therefore, our conclusions rest on the pattern rather
than on individual values. The application runs longer than \texttt{sleep}
predicts in \ExcLongerFive of \ExcN configurations, and the difference has the
same sign across backends, ratios, sizes, and placements, which run-to-run
variation would not produce systematically.

\paragraph{Future work.}
\label{sec:future}
We see two kinds of future work. The first strengthens the measurements in
this paper:
\begin{itemize}
  \item \textbf{Repeated runs.} Repeating each configuration would provide
        confidence intervals for every value and settle the single-run
        anomalies we note, such as ADIOS2 under different placements.
  \item \textbf{Idle-core effects.} Measuring core frequencies and power
        states would identify why synchronous I/O runs slower after a
        \texttt{sleep} phase.
  \item \textbf{Multi-node systems and larger outputs.} Running on a
        multi-node system with a parallel file system would expose
        contention for the network links that MPI and file system traffic
        share, which a single node cannot. In addition, outputs at the
        terabyte scale of production checkpoints would show whether the
        effect continues to grow with the output size.
  \item \textbf{PDC with more servers and larger outputs.} We measured PDC
        with one server, in the main, ratio, size, and GPU experiments, but
        not at 47\,GiB per step or at ratios above 1.25, because each run
        writes through a single server. Varying the number of servers,
        which would also make these runs affordable, and finding the cause
        of the hang with the halo exchange would complete its picture. In
        addition, a blocking or event-driven progress thread in the PDC client would likely remove most of the
        interference we measured.
\end{itemize}

The second kind opens new research directions:
\begin{itemize}
  \item \textbf{Collective, independent, and single-rank I/O.} The
        interference may depend on how the data is written. Comparing
        collective I/O, independent I/O, and I/O funneled through a single
        rank would show which access strategy interferes least with each
        type of computation.
  \item \textbf{Sensitivity of data management systems to the computational
        kernel.} Our results suggest that each data management system (DMS)
        is sensitive to a different resource. For example, the HDF5
        asynchronous VOL slowed down every working surrogate, while POSIX
        AIO mainly affected communicating ones. A systematic study of this
        sensitivity across compute-bound, memory-bound, and
        communication-bound kernels would characterize each DMS by the
        workloads it suits.
  \item \textbf{GPU data movement.} In our GPU experiments, the output is
        copied from the device to the host before it is written. The cost
        of this copy, and its interference with the background I/O, may
        depend on whether the GPU kernel is compute-bound or memory-bound,
        and on whether the data is moved to or from a running GPU kernel.
        In addition, GPUDirect Storage (GDS) writes directly from GPU memory
        to storage, which removes the host copy but may introduce its own
        interference with the GPU kernel and with the PCIe traffic of the
        application.
\end{itemize}

These directions lead to a practical question that clients of HPC systems
ask, and to a research question that follows from it.
\begin{itemize}
  \item \textbf{Client question.} \emph{My workload is network-bound,
        compute-bound, or memory-bound, and I want to overlap its I/O.
        Which DMS will give me the largest speedup?} Our results show that
        the answer depends on the workload: for a communication-heavy code,
        POSIX AIO lost less benefit than the HDF5 asynchronous VOL and
        ADIOS2, while non-blocking MPI-IO required a reserved core to help
        at all.
  \item \textbf{Research question.} \emph{Can a DMS take the
        characteristics of the client workload into account to improve the
        speedup of asynchronous I/O?} For example, a DMS could schedule
        background I/O when the computation is not using the resource the
        I/O needs, select its threading and aggregation strategy based on
        whether the workload is network-bound, compute-bound, or
        memory-bound, or reserve hardware threads only when the expected
        gain exceeds the loss of compute resources.
\end{itemize}
